\chapter{Choosing a Placement under a Training-Memory Budget}
\label{ch:budget}

The last part of the project took up a question that follows from Chapter~\ref{ch:placement}: given a limit on
training memory, which LoRA placement should be chosen? The limit is on \emph{training} memory because that is what
placement changes; inference memory is almost independent of it (Chapter~\ref{ch:deploy}). The study has two halves
with different outcomes. Following our supervisor's suggestion, the approach was built on ViT-B and Swin-B and
tested on models not used before: DeiT-B, which has the shape of ViT-B, and ViT-L, which does not.

\section{Predicting peak training memory}

Three measured properties make memory predictable. It is \emph{deterministic}: identical in all seeds, and the same
with a trained adapter as with a random one. It is \emph{linear in the batch size}: between batch sizes 8 and 256
the measured peaks deviate from a straight line by at most 0.005~GiB on ViT-B and 0.041~GiB on Swin-B. And it is set
mainly by the first adapted block (Chapter~\ref{ch:placement}).

Let $s_b$ be the size of the input of block $b$ for one image (tokens $\times$ width $\times$ 4 bytes) and $k$ the
first adapted block. Following Section~\ref{sec:bg_memory}, we model the activations stored per image as
\begin{equation}
  a \;=\; c_\text{reach} \sum_{b \ge k} s_b \;+\; c_\text{attn} \!\!\sum_{b \,\in\, \text{attn-adapted}}\!\! s_b
        \;+\; c_\text{mlp} \!\!\sum_{b \,\in\, \text{MLP-adapted}}\!\! s_b \;+\; h ,
  \label{eq:mem}
\end{equation}
and the peak for a batch of $B$ images as $\text{resident} + B \cdot a$. The first term is what every block from $k$
onwards stores because the backward pass goes through it; the second and third are the inputs that the adapters
themselves need; $h$ is a small constant for the layers outside the blocks. The resident memory (weights,
gradients, optimiser state and the input batch) follows from the parameter counts. Fitted to the seven ViT-B LoRA
configurations of Chapters~\ref{ch:lora_ft} and~\ref{ch:placement}, the constants are $c_\text{reach} = 9.95$,
$c_\text{attn} = 1.34$ and $c_\text{mlp} = 4.86$: a block that the backward pass goes through stores about ten times
the size of its input, attention adapters add about one more and MLP adapters about five more. When one ViT-B
configuration is left out of the fit, the formula predicts it to within 3.8\%.

\begin{figure}[!htb]
\centering
% generated by scripts/make_figures.py; do not edit
\begin{tikzpicture}
\begin{axis}[width=\figWone, height=\figWone, repaxis, xmin=0, xmax=22, ymin=0, ymax=22, xlabel={measured peak (GiB)}, ylabel={predicted before measuring (GiB)}, grid=major, legend style={font=\figLegendFont, draw=none, fill=none, legend cell align=left}, legend pos=north west]
\addplot[cGrey, domain=0:22, samples=2] {x};
\addlegendentry{predicted = measured}
\addplot[cAccent, only marks, mark=square*, mark size=2.2pt] coordinates {(20.287,20.12) (14.436,14.48) (8.003,7.96) (18.830,18.55)};
\addlegendentry{ViT-L}
\addplot[cAttn, only marks, mark=*, mark size=2.2pt] coordinates {(5.459,5.51) (5.059,5.11) (1.462,1.44) (4.660,4.70) (4.260,4.29) (3.860,3.88) (3.461,3.48) (3.061,3.07) (2.262,2.25)};
\addlegendentry{DeiT-B}
\addplot[cBoth, only marks, mark=triangle*, mark size=2.2pt] coordinates {(2.560,2.52) (7.175,6.85) (1.815,0.90) (6.000,5.76) (4.854,4.68)};
\addlegendentry{Swin-B}
\node[font=\figLegendFont, anchor=west, align=left] (miss) at (axis cs:6.5,2.90) {Swin-B, start block 22:\\predicted 0.90,\\measured 1.82};
\draw[->, cGrey] (miss.west) -- (axis cs:2.17,0.95);
\end{axis}
\end{tikzpicture}

\caption{Peak training memory at batch size 64 for 18 configurations, predicted with the ViT-B constants before
measuring, against the measurement. ViT-L has a new width and depth; DeiT-B has the shape of ViT-B; Swin-B: start
blocks 2 to 22.}
\label{fig:mem_pred}
\end{figure}

\paragraph{Test.} Before measuring them, we wrote down the predicted peak of 18 new configurations: attention LoRA
on every block from a start block $k$ to the last on DeiT-B ($k = 0$ to 10) and Swin-B ($k = 2$ to 22), and four
configurations of ViT-L (Figure~\ref{fig:mem_pred}).
\begin{itemize}
  \item \textbf{ViT-L} is the real test, because its width and depth differ from those of ViT-B: the mean error is
        0.8\% and the largest 1.5\%. The formula also predicted that three of its configurations would not fit on
        the partition at batch size 256, and all three ran out of memory.
  \item \textbf{DeiT-B} has the shape of ViT-B, so its agreement (within 1.7\%) shows that memory depends on the
        shape and not on the weights; it is a consistency check. Each block by which the start moves later saves
        0.40~GiB, at every step (Figure~\ref{fig:mem_start}).
  \item \textbf{Swin-B} is predicted to within 4.6\% for start blocks 2 to 16, by constants fitted on a different
        architecture (1.8\% with constants fitted on Swin-B itself).
  \item \textbf{One prediction failed.} For Swin-B with adapters on the last two blocks only ($k=22$) the formula
        predicts 0.90~GiB and the measurement is 1.82~GiB. The formula counts activations that are stored until the
        backward pass, which are small here (0.27~GiB); the peak is instead set by memory used briefly during the
        forward pass through the large early stages. The formula needs a lower bound for this temporary memory,
        which we have not built. Tight budgets lead exactly to such late starts, so this limit matters.
\end{itemize}

On the memory side, the budget problem therefore needs no search: the budget fixes the earliest possible start
block. A budget of 4.5~GiB at batch size 64, for example, is met on ViT-B and DeiT-B by starting at block 3
(4.26~GiB). Which of the placements that fit is the most accurate is a separate question.

\section{Predicting accuracy: a rule that failed}

For the memory budget, the question is how much accuracy is lost when the adapters start at block $k$ instead of
block 0. Before training any such configuration, we wrote down a simple rule for this loss: it grows in proportion to
the fraction of the blocks that are skipped,
\[
  \Delta\text{acc}(k) \;\approx\; -\beta\,\frac{k}{L},
\]
where $\Delta\text{acc}$ is the change in test accuracy against adapters on all blocks, $k$ the first adapted block
and $L$ the number of blocks. The constant $\beta$ came from the two second-half results of
Chapter~\ref{ch:placement}, which start half-way through the model ($k/L = 0.5$): a loss of 1.56~pp on ViT-B and
1.18~pp on Swin-B gives $\beta = 3.1$ and $2.4$~pp, on average 2.7~pp. For a start at block 3 of 12 ($k/L = 0.25$)
the rule predicts a loss of 0.7~pp on DeiT-B (average $\beta$) and 0.8~pp on ViT-B (its own $\beta$), and for a start
at block 6 on DeiT-B a loss of 1.4~pp. Each prediction was given a range of about $\pm$0.6~pp, chosen by judgement,
not computed. The rule predicts a change, not an accuracy, because accuracy levels differ between models (DeiT-B is
about 4.5~pp below ViT-B) and because a change, compared seed by seed with the all-blocks run, is far less affected
by the variation between seeds and between sets of runs. Table~\ref{tab:acc_start} compares the predictions with the
measurements.

\begin{table}[!htb]
\centering\footnotesize
\setlength{\tabcolsep}{3.5pt}
\begin{tabular}{lccccc}
\toprule
attention LoRA & all blocks & blocks $k$ to last & predicted & measured & within the \\
 & (\%) & (\%) & change (pp) & change (pp) & prediction? \\
\midrule
DeiT-B, $k=3$ & 86.56 & 86.71 & $-$0.7 ($-$1.4 to $-$0.2) & $+$0.15 $\pm$ 0.04 & no \\
DeiT-B, $k=6$ & 86.56 & 86.24 & $-$1.4 ($-$2.5 to $-$0.7) & $-$0.32 $\pm$ 0.35 & no \\
ViT-B, $k=3$  & 91.28 & 90.97 & $-$0.8 ($-$1.4 to $-$0.2) & $-$0.31 $\pm$ 0.27 & yes, at the edge \\
\bottomrule
\end{tabular}
\caption{Accuracy of attention LoRA on all blocks and on blocks $k$ to the last (mean over 3 seeds), and the change
between them: predicted with the rule above before the runs, and measured (mean $\pm$ standard deviation of the
per-seed differences; negative = less accurate). Details: Table~\ref{tab:acc_pred} in the appendix.}
\label{tab:acc_start}
\end{table}

The rule did not hold: a later start cost less than the straight line predicted. On DeiT-B, starting at block 3
cost no accuracy in these runs, with 22\% less training memory; on ViT-B it cost 0.31~pp, not much more than the
0.2~pp by which two sets of ViT-B runs with the same seeds can differ (Appendix~\ref{app:impl}). The memory and time predictions for the same
runs held: the training runs reproduced the profiled peak memory to four decimals (4.2600 and 3.0610~GiB), and the
training epoch and the backward pass were predicted to within 1~s and 1~ms. The accuracy of a placement therefore
cannot be read from the shape of the model in the way its memory can, and we did not find another way to predict it.
